% MT020 / explanation — canonical editable content.
模型连训练样本都处理不好，说明它还没有充分表示训练数据中的关系。这叫欠拟合（underfitting）。本题已经给了大量训练样本，重点应放在“线性规则是否足够表达症状之间的关系”。


线性分类器的分数是\(w_1x_1+\cdots+w_dx_d+b\)：每个症状各自贡献一项，再相加。每个症状的作用若取决于另一个症状是否出现，仅把各项相加就可能不够。补充例子：两项特征各取0或1，只有恰好一项为1时是正类，即\((1,0)\)、\((0,1)\)为正，\((0,0)\)、\((1,1)\)为负。这四个角无法由一条直线分开。不断增加相同形式的数据也不会让直线获得弯曲能力。


可以提出两种不同的办法：



\begin{enumerate}[start=1]

\item \textbf{保留线性模型，改善输入表示。} 加入有依据的交互项，如\(x_1x_2\)，或多项式特征。模型对新特征仍是线性的，但在原输入空间能表示更复杂的关系。对上述例子，加入\(z=x_1x_2\)，取分数\(f=x_1+x_2-2z-0.5\)：两个正类点的分数都是0.5，两个负类点都是\(-0.5\)，按分数正负就能分开。

\item \textbf{换成能表达非线性关系的模型。} 例如核SVM或树模型，再用验证集比较是否真正改善。

\end{enumerate}


若正则化太强，减弱惩罚也可能改善欠拟合。实际排查还要确认训练程序已收敛、标签没有系统错误；按本题给定条件，主要诊断与两项答案仍是表示能力不足及上述改进。


\textbf{出处：} 2021A期中 Q7，原卷第3页（10分）


\textbf{答案依据：} 2021A期中 Q7 随卷／配套参考答案；条件与解释经复核。
